\section{観測境界と脅威モデル}

\subsection{観測境界}

対象は SCSI/NVMe 的なブロックストレージである\cite{nvme}。検出器は装置内、
または装置に近いストレージ隣接プロセッサ上に配置され、表
\ref{tab:observable} の情報を集約する。

\begin{table}[tb]
  \centering
  \caption{観測可能情報と採用方針}
  \label{tab:observable}
  \begin{tabularx}{\linewidth}{lllY}
    \toprule
    情報 & 必須 & 収集コスト & 採用方針 \\
    \midrule
    時刻 & yes & 低 & 10 秒窓への集約に用いる。 \\
    read/write 種別 & yes & 低 & write ratio、カウント、バイト数に使う。 \\
    LBA/offset & yes & 低 & 10 秒窓の平均 LBA に集約する。 \\
    転送長 & yes & 低 & 10 秒窓の平均転送長に集約する。 \\
    窓間変化 & optional & 低 & 前フレームとの差分だけを保持する。 \\
    圧縮率/entropy 相当 & optional & 中 & 既存の装置カウンタがある場合だけ使う。 \\
    ペイロード走査 & no & 高 & 初期スコープから除外する。 \\
    ファイル名/プロセス & no & 不可 & ブロック装置単体では観測しない。 \\
    \bottomrule
  \end{tabularx}
\end{table}

この境界は、エンドポイント型検出に比べて意図的に情報量が少ない。
その代わり、ホスト OS、プロセス、ファイルシステムの信頼に依存しない。
この選択は、強い意味でのランサムウェア分類ではなく、
``ストレージから見た異常な暗号化・大量書き換え挙動の検出'' を目標にする。

\subsection{脅威モデル}

攻撃者は、保護対象ブロックデバイスへ正規にアクセスできるホスト上で
ランサムウェアを実行する。ストレージ装置は、プロセス名、ファイル拡張子、
ユーザ意図、ファイル内容の意味を知らない。検出器が期待する攻撃時の
観測変化は、write-heavy burst、write ratio の変化、平均 I/O 長の変化、
平均 LBA の移動、圧縮率低下または entropy 上昇である。実際の上書き量や
ファイル単位の暗号化進行は装置内 scalar から直接は観測せず、評価時に
raw trace と攻撃ラベルから推定する。

表 \ref{tab:anomaly-taxonomy} は、本稿で扱う異常をストレージから見える信号、
混同しやすい正常処理、必須評価に分解したものである。この分類により、
本研究の対象は malware family の識別ではなく、ブロック装置から観測できる
ransomware-like storage anomaly の検出であることを明確にする。

\begin{table}[tb]
  \centering
  \caption{ストレージ観測に基づくランサムウェア異常分類}
  \label{tab:anomaly-taxonomy}
  \small
  \begin{tabularx}{\linewidth}{lYYY}
    \toprule
    異常型 & 観測信号 & 主な confounder & 必須評価 \\
    \midrule
    write burst &
    write count/bytes の急増、write ratio の変化 &
    backup、database checkpoint、bulk copy &
    write-throughput baseline と false alarms per volume per day。 \\
    LBA locality shift &
    平均 LBA と窓間移動量 &
    file-system aging、仮想ディスク copy-on-write、defrag &
    LBA ablation、volume condition holdout。 \\
    I/O length shift &
    平均転送長と窓間変化 &
    database、log compaction、snapshot merge &
    length ablation、workload-family holdout。 \\
    overwrite impact &
    write burst、平均 LBA/長さの窓間変化 &
    restore、scrub、deduplication、compression job &
    raw trace から測る time to detect と alert 前の overwritten bytes/LBA fraction。 \\
    payload randomness proxy &
    圧縮率低下、entropy 相当信号上昇 &
    暗号化済み volume、圧縮済み benign file、高 entropy file &
    metadata-only variant と高 entropy benign replay。 \\
    \bottomrule
  \end{tabularx}
\end{table}

一方で、低速に暗号化するランサムウェア、少数ファイルだけを狙う攻撃、
既に暗号化されたボリューム、バックアップ、圧縮、重複排除、データベース
メンテナンス、仮想ディスクの copy-on-write などは、検出困難または
false positive の主要因になり得る。このため、検出器は単一の anomaly score
だけでなく、特徴チャネルごとの誤差寄与を出力する必要がある。
ただし、この誤差寄与は ``どの scalar 統計が外れたか'' を示すものであり、
どのファイルやどの LBA が暗号化されたかを復元するものではない。
